%% CAPITULO 2 -- REFERENCIAL TEORICO

\section{Referencial teórico}

A triagem de manifestações consiste em identificar o assunto de cada demanda e atribuí-la à unidade competente. Na literatura recente, essa tarefa aparece tratada como classificação de textos, com modelos treinados sobre demandas reais de plataformas de participação e de agências públicas \parencite{buryakov2024,majumder2025}.

\citeonline{buryakov2024} propõem um classificador multirrótulo, baseado em um modelo BERT ajustado, para atribuir petições online aos departamentos governamentais responsáveis, treinado com dados de uma plataforma de petições de Taiwan. A Tabela~\ref{tab:estudos-base} resume o contexto e a tarefa dos três estudos de referência adotados neste esqueleto.

\begin{table}[htb]
  \caption{Estudos de referência sobre triagem de demandas públicas}
  \label{tab:estudos-base}
  \centering
  \begin{tabular}{p{3.0cm}p{5.0cm}p{5.0cm}}
    \toprule
    \textbf{Obra} & \textbf{Contexto institucional} & \textbf{Tarefa tratada} \\
    \midrule
    \citeonline{buryakov2024} & Plataforma de petições de Taiwan & Atribuição de petições a departamentos \\
    \citeonline{majumder2025} & Agência ambiental dos Estados Unidos & Filtragem e roteamento de denúncias \\
    \citeonline{prencipe2026} & Administração pública regional na Itália & Preparo de respostas a consultas, com humano no circuito \\
    \bottomrule
  \end{tabular}
  \fonte{Elaboração própria, com base nos resumos das obras citadas.}
\end{table}

Ao descrever o problema enfrentado por agências reguladoras que dependem de denúncias do público, \citeonline{majumder2025} afirmam, no resumo do trabalho:

\begin{citacao}
Regulatory agencies often operate with limited resources and rely on tips from the public to identify potential violations. However, processing these tips at scale presents significant operational challenges, as agencies must correctly identify and route relevant tips to the appropriate enforcement divisions. \parencite{majumder2025}
\end{citacao}

No contexto de fluxos de trabalho administrativos, \citeonline{prencipe2026} apresentam um sistema multiagente com humano no circuito, que combina agentes especializados e recuperação híbrida sobre um grande acervo de documentos administrativos, e relatam seis meses de uso em produção.
